\documentclass[12pt]{article}
\usepackage[margin=1in]{geometry}
\usepackage{enumitem}
\usepackage{titlesec}
\usepackage{parskip}
\usepackage{graphicx}
\usepackage{xcolor}
\usepackage{hyperref}
\usepackage{fontenc}

\title{\textbf{Nichter 2008 RG Notes}\\
\large Nichter, Simeon. ``Vote Buying or Turnout Buying? Machine Politics and the Secret Ballot'' [2008].\\
\textit{American Political Science Review} 102(1): 19-31.}
\author{}
\date{}

\begin{document}

\maketitle

Add Stokes 2005 to reading list

Core puzzle: how does vote-buying sustain itself when political actors cannot perfectly monitor voter behavior (who they cast their vote for)? Note that Stokes et al 2013 largely ``ignores this away'' and suggests that the monitoring gap is actually quite small; here, Nichter provides an alternative explanation that allows for a more significant monitoring gap

Provides both a theoretical model of turnout-buying and an empirical comparison of Stokes 2005's data instead specifically analyzing turnout buying

Rests on the (virtually uncontested) assumption that it is easier to monitor turnout than specific individual votes

Categorize voters not just by their preference of candidates, but also by their likelihood to vote. (Perhaps a broader way to model a three-way ranking of preferences?)

Turnout-buying requires some influence of future interactions: there must be some reason for the payee to not shirk, so accordingly we assume there is some consideration of repetition (shared assumption w/ Stokes et al)

Model predicts that turnout buying should be targeted towards those already aligned with the machine

\end{document}
